\chapter*{DATOS GENERALES}
\addcontentsline{toc}{chapter}{DATOS GENERALES}
\vspace{0.35cm}

\begingroup
\singlespacing
\setlength{\parskip}{0pt}
\setlength{\parindent}{0pt}
\setlength{\tabcolsep}{10pt}
\renewcommand{\arraystretch}{1.28}
\noindent
\begin{tabularx}{\textwidth}{
  >{\bfseries\RaggedRight\setlength{\parindent}{0pt}\arraybackslash}p{4.15cm}
  J
}
\toprule
\rowcolor{lightgray}
\textbf{Elemento} & \textbf{Descripción} \\
\midrule
Título del plan & \tituloTesis. \\
\addlinespace[5pt]
Tesistas & \mbox{Luis Enrique Koc Góngora};\newline
           \mbox{Herbert Antonio Meléndez García}. \\
\addlinespace[5pt]
Asesor & Samuel Oporto. \\
\addlinespace[5pt]
Programa & Maestría en Inteligencia Artificial. \\
\addlinespace[5pt]
Unidad académica & Unidad de Posgrado de la Facultad de Ingeniería Industrial y de Sistemas, Universidad Nacional de Ingeniería. \\
\addlinespace[5pt]
Área involucrada & Inteligencia artificial científica, modelado computacional y sistemas eléctricos de potencia. \\
\addlinespace[5pt]
Lugar de desarrollo & Entorno computacional de la Facultad de Ingeniería Industrial y de Sistemas de la Universidad Nacional de Ingeniería (FIIS-UNI) y casos de referencia documentados de instalaciones de cables eléctricos enterrados. \\
\addlinespace[5pt]
Duración estimada & Doce meses, desde julio de 2026 hasta junio de 2027. \\
\addlinespace[5pt]
Fecha de inicio & Julio de 2026. \\
\addlinespace[5pt]
Fecha de culminación prevista & Junio de 2027. \\
\addlinespace[5pt]
Tipo de graduación & Maestría de especialización o profesionalizante. \\
\addlinespace[5pt]
Artefacto de\newline investigación & Modelo y prototipo computacional 2D basado en redes neuronales informadas por física para analizar campos térmicos y ampacidad en entornos heterogéneos. \\
\bottomrule
\end{tabularx}
\endgroup
